% Generated draft: requires mathematical and source review.
\documentclass[11pt]{article}
\usepackage[margin=25mm]{geometry}
\usepackage{fontspec}
\setmainfont{Latin Modern Roman}
\usepackage{amsmath,amssymb,mathtools,unicode-math}
\setmathfont{Latin Modern Math}
\usepackage{enumitem,needspace,xurl,hyperref}
\setlength{\parindent}{0pt}
\setlength{\parskip}{5pt}
\setlength{\emergencystretch}{3em}
\newcommand{\N}{\mathbb{N}}
\newcommand{\Z}{\mathbb{Z}}
\newcommand{\Q}{\mathbb{Q}}
\newcommand{\R}{\mathbb{R}}
\newcommand{\C}{\mathbb{C}}
\newcommand{\F}{\mathbb{F}}
\newcommand{\A}{\mathbb{A}}
\newcommand{\PP}{\mathbb{P}}
\newcommand{\E}{\mathbb{E}}
\newcommand{\eps}{\varepsilon}
\newcommand{\dd}{\,\mathrm{d}}
\begin{document}
\section*{UnsolvedMath ID: 800625}
\textbf{Source problem code:} catalogue entry
\subsection{Definitions and mathematical statement}
Equip \(\R\) with its usual metric, topology, and addition, and let \(\N=\{1,2,\ldots\}\). A set \(X\subseteq\R\) has \emph{strong measure zero} if, for every sequence \((\varepsilon_n)_{n\in\N}\) of positive real numbers, there are open intervals \(I_n\subseteq\R\) such that
\[
X\subseteq\bigcup_{n\in\N} I_n
\quad\text{and}\quad
\operatorname{length}(I_n)<\varepsilon_n\quad(n\in\N).
\]
A set is \emph{meager} if it is a countable union of nowhere dense sets, where a set is nowhere dense when its closure has empty interior. A set \(X\subseteq\R\) is \emph{meager-additive} if, for every meager set \(M\subseteq\R\), the sum
\[
X+M=\{x+m:x\in X,\ m\in M\}
\]
is meager in \(\R\). \emph{Borel's Conjecture} asserts that every strong measure zero subset of \(\R\) is at most countable.

The supplied problem asks whether the following implication holds:
\[
\begin{gathered}
\bigl[\forall X\subseteq\R,\quad
X\text{ has strong measure zero}
\Longrightarrow X\text{ is meager-additive}\bigr]
\\[2pt]
\Longrightarrow
\bigl[\forall X\subseteq\R,\quad
X\text{ has strong measure zero}
\Longrightarrow |X|\leq\aleph_0\bigr].
\end{gathered}
\]
Under the usual foundational interpretation, the question is whether this implication is provable in ZFC. The supplied record does not explicitly specify a foundational theory; this interpretation requires confirmation against the original source.
\subsection{Short English statement}
If every strong measure zero set of reals remains meager after addition to any meager set, must every strong measure zero set of reals be countable? This is the supplied question asking whether that hypothesis implies Borel's Conjecture.
\subsection{Sources}
\begin{itemize}
\item Local UnsolvedMath registry, permanent ID \(800625\), source problem code \texttt{catalogue entry}; supplied title: Selection Principles in Mathematics --- Problem 54 ([BarJu93]).
\item Tsaban, \emph{Selection Principles and Special Sets of Reals: Open Problems} (2006), Problem 54. Supplied source URL: \url{https://arxiv.org/abs/math/0609601}. The record reports access on 2026-10-16 and extraction method \texttt{arxiv\_latex}; these are supplied provenance details, not independently verified access.
\item The supplied attribution/context is the reference label [BarJu93]. No full bibliographic entry for that label is included in the record.
\item The registry attributes the problem to Tsaban, gives proposed year 2006, and labels its status as open. This status is supplied curation, not an independently verified literature assessment.
\end{itemize}
\end{document}
